%! AP Statistics — Unit 1, Lesson 1.8 — Student Packet Cover Sheet
\documentclass[10pt]{article}
\usepackage{apstats-article}
\usepackage{apstats-boxes}

%=============================================================================
\begin{document}

% ── Full-bleed navy banner ────────────────────────────────────────────────────
\begin{tikzpicture}[remember picture, overlay]
  \fill[navy] (current page.north west)
    rectangle ([yshift=-0.9in]current page.north east);
\end{tikzpicture}
\vspace{-0.8in}

\begin{center}
  {\color{white}\textbf{\LARGE AP Statistics}}\\[4pt]
  {\color{skymid}\large\textbf{Unit 1: Exploring One-Variable Data}}\\[6pt]
  {\color{white}\textbf{\large Lesson 1.8 \quad Graphical Representations of Summary Statistics}}
\end{center}

\vspace{0.22in}

\namedateperiod

\vspace{0.18in}

% ── LEARNING TARGETS ─────────────────────────────────────────────────────────
\begin{learningtargetbox}
\small
\begin{enumerate}[leftmargin=1.5em, itemsep=5pt]
  \item \ldots{} read a boxplot to identify the five-number summary
        ($\min, Q_1, M, Q_3, \max$) and estimate the IQR and range.
  \item \ldots{} determine whether a value is an outlier using the
        \textbf{$1.5 \times \text{IQR}$ fence rule}, and explain the result in context.
  \item \ldots{} interpret a \textit{modified boxplot}, distinguishing the
        whisker endpoint from the fence and from plotted outlier dots.
  \item \ldots{} describe a distribution shown by a boxplot using \textbf{SOCS}
        and explain what a boxplot cannot show compared to a histogram.
\end{enumerate}
\end{learningtargetbox}

\vspace{0.14in}

% ── TABLE OF CONTENTS ────────────────────────────────────────────────────────
\begin{tocbox}
\small
\renewcommand{\arraystretch}{1.5}
\begin{tabularx}{\linewidth}{@{} c l X r @{}}
\rowcolor{sky}
\textbf{\#} & \textbf{Component} & \textbf{Description} & \textbf{Score} \\
\midrule
1 & Warm-Up        & Read and interpret a pre-drawn boxplot; apply L1.7 skills & \blank{1.2cm} \\
2 & Guided Notes   & Boxplot anatomy, modified boxplots, outliers, SOCS         & \blank{1.2cm} \\
3 & Group Activity & Compute five-number summary; interpret a modified boxplot   & \blank{1.2cm} \\
4 & Exit Ticket    & Read a modified boxplot; apply the outlier rule; SOCS      & \blank{1.2cm} \\
5 & Homework       & Multi-part practice: interpret, outlier check, AP-style    & \blank{1.2cm} \\
\midrule
  & \multicolumn{2}{r}{\textbf{Total}} & \blank{1.2cm} \\
\end{tabularx}
\end{tocbox}

\vspace{0.14in}

% ── KEEP IN MIND ─────────────────────────────────────────────────────────────
\begin{remindbox}
\small
In Lesson 1.7 you calculated the five-number summary and the IQR fence rule.
Lesson 1.8 puts those numbers on a number line as a \textit{boxplot} --- a
visual summary of spread, center, and outliers in one compact display.

\vspace{6pt}
\begin{tabularx}{\linewidth}{@{} >{\centering\arraybackslash\columncolor{goldbg}}p{0.06\linewidth}
                                    X @{}}
\textbf{\large!} &
\textbf{Each section of a boxplot contains roughly 25\% of the data.}\enspace
The box spans $Q_1$ to $Q_3$ (the middle 50\%); each whisker covers one more
quarter. This is an approximation --- actual percentages depend on the data. \\[6pt]
\textbf{\large!} &
\textbf{The whisker ends at the last \textit{non-outlier} value, not at the fence.}\enspace
In a modified boxplot, the fence ($Q_3 + 1.5\,\text{IQR}$) is an invisible
threshold. The whisker stops at the actual data point just below the fence. \\[6pt]
\textbf{\large!} &
\textbf{Boxplots hide shape, gaps, and multimodality.}\enspace
A dataset with two distinct clusters can produce the same boxplot as a
unimodal dataset. Always pair a boxplot with context and, when possible,
a histogram or dotplot. \\[6pt]
\textbf{\large!} &
\textbf{Outliers must be investigated, not deleted.}\enspace
A value flagged by the $1.5 \times \text{IQR}$ rule may be a data-entry error
\textit{or} a genuine extreme --- check the source before reporting. \\
\end{tabularx}
\end{remindbox}

\end{document}
